\section{Experimental protocol}
\label{sec:protocol}

The experiments use all twenty instances in Set~A. Their sizes range from 20 to 400 vertices, 80 to
2000 directed links, and 20 to 6000 demands. Two experiments are performed for each instance. The
first considers only period~0 and omits the reconfiguration budget, as required in Step~6 of the
project. The second considers periods~0 and~1 together and includes the budget constraint, as
required in Steps~7--9.

The period-0 time limit is 1800 seconds (30 minutes), while the two-period limit is 3600 seconds (one
hour). Each experiment is repeated ten times using seeds 2026--2035. Every solver process uses one
thread. The code is run with Julia~1.13.2, JuMP~1.32.1~\citep{jump2017}, and HiGHS~1.26.0 on 64-bit
Linux~6.8.0. The machine has two Intel Xeon Platinum 8558 processors, with 44 physical cores per
socket and two hardware threads per core (88 physical cores and 176 logical CPUs in total), and
1.7\,TiB of RAM. Although the machine has many cores, the experiments are executed sequentially: only
one instance and one seed are active at a time, and HiGHS uses one thread. Thus, no two experimental
jobs compete for CPU or memory. The operating-system background load was not logged, so no numerical
load average is claimed.

For every run, the output file records the instance size, checked MLU, best valid lower bound,
relative gap, solver time, number of branch-and-bound nodes, solution status, and solution method. In
accordance with the project instructions, \emph{solver time} includes only elapsed time inside calls
to \texttt{optimize!}. When an algorithm makes several solver calls, these times and their
branch-and-bound node counts are summed. Data loading, graph construction, coefficient computation,
pricing, validation, and file writing are excluded. A run may stop before its limit when the lower
bound certifies the incumbent. Small overruns are possible because a solver call is timed when it
returns.

The original experiment files do not contain end-to-end wall-clock time. Consequently, the tables
report the required solver time but do not present it as the total running time of the algorithm. In
particular, a small solver time for a column-generation run does not include shortest-path
preprocessing, pricing, or column construction. Measuring total wall-clock time would require
rerunning the experiments with an outer timer.

Tables~\ref{tab:t0} and~\ref{tab:full} report the best MLU, mean MLU, sample standard deviation of
MLU, mean solver time, mean gap, mean node count, and number of certified runs out of ten. The gap of
one run is
\[
\mathrm{gap}=\max\left\{0,\,1-\frac{\mathrm{LB}}{\mathrm{UB}}\right\},
\]
where UB is the checked MLU and LB is a valid lower bound. The \emph{Cert.} column makes mixed outcomes
explicit; for example, 8/10 means that eight runs certify the scalar MLU optimum and two do not. An
infeasible run would be reported separately, but every run in this study produced a feasible
incumbent.
